%\documentclass[12pt,a4paper]{article}
\documentclass[german,bibnum,beleg,zihtitle,german,hyperref,utf8]{zihpub}
\author{Maximilian Knespel}
\betreuer{Dipl.-Inf. Nico Hoffmann}
\title{Belegarbeit Computational Science and Engineering\\Verteilte GPGPU-Berechnungen mit Spark }
\date{}
\bibfiles{literatur}

\usepackage[german,onelanguage]{algorithm2e}

\usepackage[utf8]{inputenc}
\usepackage[T1]{fontenc}
	% if you don't use \usepackage[T1]{fontenc},
    % - Words containing accented characters (ö,...) cannot be automatically hyphenated,
    % - You cannot properly copy-and-paste srom the output (DVI/PS/PDF),
    % - Characters like the pipe sign, less than and greater sign give unexpected results in text.
\usepackage{lmodern}  % Vector Fonts even when using T1 !
%\usepackage[ngerman]{babel}
\usepackage[dvipsnames,cmyk,table]{xcolor} % xcolor and more colornames, cymk for printing
\usepackage{cmap} % pro­vides char­ac­ter map ta­bles, which makes PDFs search­able and copy-able (I don't see a difference yet!!!)
\usepackage[labelfont=bf,format=plain]{caption} % allows captions in non-figure environments with \captionof and \captionsetup
\usepackage[binary-units=true]{siunitx}    % SI units paper conform with \SI{3}{\nano\seconds}
\usepackage{tabularx}
\newcolumntype{Y}{>{\centering\arraybackslash}X}

\definecolor{mygreen}{rgb}{0,0.6,0}
\definecolor{mygray}{rgb}{0.5,0.5,0.5}
\definecolor{mymauve}{rgb}{0.58,0,0.82}

\usepackage{listings}   % source code listings
\usepackage{textcomp}
\usepackage{courier}
\usepackage{float}      % enables \begin{figure}[H] to disable floating figure
\usepackage{ulem}

\lstset{ %
  backgroundcolor=\color{white},   % choose the background color; you must add \usepackage{color} or \usepackage{xcolor}
  basicstyle=\footnotesize\ttfamily, % the size of the fonts that are used for the code
  breakatwhitespace=false,         % sets if automatic breaks should only happen at whitespace
  breaklines=true,                 % sets automatic line breaking
  captionpos=b,                    % sets the caption-position to bottom
  commentstyle=\color{mygreen},    % comment style
%  deletekeywords={...},            % if you want to delete keywords from the given language
  escapeinside={\%*}{*)},          % if you want to add LaTeX within your code
  extendedchars=false,              % lets you use non-ASCII characters; for 8-bits encodings only, does not work with UTF-8
  %frame=single,                    % adds a frame around the code
  keepspaces=true,                 % keeps spaces in text, useful for keeping indentation of code (possibly needs columns=flexible)
  keywordstyle=\color{blue},       % keyword style
  morekeywords={*,...},            % if you want to add more keywords to the set
  numbers=left,                    % where to put the line-numbers; possible values are (none, left, right)
  numbersep=5pt,                   % how far the line-numbers are from the code
  numberstyle=\tiny\color{mygray}, % the style that is used for the line-numbers
  rulecolor=\color{black},         % if not set, the frame-color may be changed on line-breaks within not-black text (e.g. comments (green here))
  showspaces=false,                % show spaces everywhere adding particular underscores; it overrides 'showstringspaces'
  showstringspaces=false,          % underline spaces within strings only
  showtabs=false,                  % show tabs within strings adding particular underscores
  stepnumber=1,                    % the step between two line-numbers. If it's 1, each line will be numbered
  stringstyle=\color{mymauve},     % string literal style
  tabsize=2,                       % sets default tabsize to 2 spaces
  title=\lstname,                  % show the filename of files included with \lstinputlisting; also try caption instead of title
  xleftmargin=.25in,               % indent listings by default
  upquote=true                     % single quotes '' appear straight instead of as backticks. Needs textcomp
}

\usepackage{amsmath, amsthm, amssymb}
\allowdisplaybreaks
\setlength\parindent{0pt} % stop indenting every paragraph by 1cm. It looks awful!
\usepackage{esint} % for oiint
\usepackage{attachfile} % insert attachments into pdf O_O

%\usepackage{pdfpages}   % include single pages from larger pdfs (this is a pdfxcviewerpro feature, it's cool that latex can do this for free!!)

\usepackage{enumitem}
\setitemize{itemsep=0pt,topsep=3pt}
\setenumerate{itemsep=0pt,topsep=3pt}

% http://tex.stackexchange.com/questions/2644/how-to-prevent-a-page-break-before-an-itemize-list
\makeatletter
\newcommand{\nolisttopbreak}{\vspace{\topsep}\nobreak\@afterheading}
\makeatother


%%%%%%%%%%%%%%%%%%%%%%%%%%%%%%%%%%%%%%%%%%%%%%%%%%%%%%%%%%%%%%%%%%%%%%%%%%%%%%%%
% main part
%%%%%%%%%%%%%%%%%%%%%%%%%%%%%%%%%%%%%%%%%%%%%%%%%%%%%%%%%%%%%%%%%%%%%%%%%%%%%%%%

\begin{document}

\hyphenchar\font=\string"7F % make "= unneeded
%\hyphenation{Kon-fi-gu-ra-ti-ons-tool} % why isn't this working damn it!?

% All files must be UTF8-without BOM encoded!!!

%sed -e 's/\\section/\\chapter/'          -i Aufgabe*.tex
%sed -e 's/\\subsection/\\section/'       -i Aufgabe*.tex
%sed -e 's/\\subsubsection/\\subsection/' -i Aufgabe*.tex

%%%%%%%%%%%%%%%%%%%%%%%%%%%%%%%%%%%%%%%%%%%%%%%%%%%%%%%%%%%%%%%%%%%%%%%%%%%%%%%%
\chapter{Einführung}
%%%%%%%%%%%%%%%%%%%%%%%%%%%%%%%%%%%%%%%%%%%%%%%%%%%%%%%%%%%%%%%%%%%%%%%%%%%%%%%%

Im Rahmen dieser Belegarbeit soll ein Ansatz entwickelt werden, um mit Java oder Scala auf heterogenen Clustersystem mit Grafikkarten zu rechnen. Es wurde sich für eine Kombination von Spark für die Kommunikation im Cluster und Rootbeer für die Grafikkartenprogrammierung entschieden.

Für Spark wurde sich entschieden, weil es nicht nur eine einfache Programmierung von Clustern mittels des MapReduce-Programmiermodells ermöglicht, sondern dazu noch zahlreiche Bibliotheken z.B. für Maschinenlernen zur Verfügung stellt.

Rootbeer wurde genommen, weil es das Schreiben von CUDA-Kernel aus Java heraus erlaubt. Die so geschriebenen Kernel können dann sowohl auf Grafikkarten als auch auf dem Host ausgeführt werden.

Zuerst wird in den Kapiteln \ref{sct:montecarloalgo}-\ref{sct:spark} die benutzten Algorithmen und Bibliotheken vorgestellt, in Kapitel~\ref{sct:implementation} wird die eigene Implementierung dokumentiert und in Kapitel~\ref{sct:benchmarks} werden Benchmarks dieser Implementierung vorgestellt.

%Dabei soll die Vorgehensweise reproduzierbar dokumentiert werden, mit dem Ziel, die Programmierung auf Grafikkarten-Clustern so einfach wie möglich zu machen. Z.B. durch ein Skript. Als ersten Testfall wird ein Monte-Carlo-Algorithmus zur Berechnung von Pi implementiert, da dieser sehr rechenlastig ist. Weiterhin soll ein kommunikationslastiger Algorithmus wie z.B. Mergesort auf einem Grafikkartencluster untersucht werden.


%%%%%%%%%%%%%%%%%%%%%%%%%%%%%%%%%%%%%%%%%%%%%%%%%%%%%%%%%%%%%%%%%%%%%%%%%%%%%%%%
\chapter{Apache Spark}
\label{sct:spark}
%%%%%%%%%%%%%%%%%%%%%%%%%%%%%%%%%%%%%%%%%%%%%%%%%%%%%%%%%%%%%%%%%%%%%%%%%%%%%%%%

Spark ist ein Programmierframework für Datenanalyse auf Clustern, was vor allemim Zusammenhang mit ''Big Data'' an Beliebtheit gewonnen hat.
Es vereinigt hierbei ausfallgehärtet Funktionalitäten von Batchverarbeitungssystemen bzw. Cluster-Management wie z.B. SLURM, Kommunikation zwischen Prozessen, wie z.B. OpenMPI und OpenMP sie zur Verfügung stellen, und zusätzliche problemspezifische Programmierbibliotheken.
Spark stellt Programmierschnittstellen für Java, Scala und Python zur Verfügung. Für Scala und Python existieren interaktive Eingabeaufforderungen mit der z.B. interaktive Echtzeitanalysen auf Clustern mit Spark möglich sind.\\

Apache Spark oder auch Spark Core stellt die Grundfunktionalität für verteiltes Rechnen bereit.
Das inkludiert die erwähnte Ausfallsicherheit, Management von Knoten über ein Web-Interface und Prozess-Scheduling\cite{learningspark}.
Die Programmierschnittstelle hierfür sind der Spark-Kontext und die Resilient Distributed Datasets (RDD).

\begin{figure}
    \centerline{\includegraphics[width=0.5\linewidth]{spark-stack.png}}
    \caption{Zusammensetzung des Spark-Frameworks\cite{spark}}
    \label{fig:sparkstack}
\end{figure}

Beispielhaft für den Spark-Stack seien hier die Bibliotheken MLlib\cite{mllib} und GraphX erwähnt.

MLlib, kurz für Machine Learning Library, umfasst Funktionen wie lineare Regression, den K-Means-Algorithmus, Latent Dirichlet Allocation, Hauptkomponentenanalyse, das für Maschinenlernen benötigte stochastische Gradientenverfahren, u.v.m.

% GraphX
GraphX erweitert Spark-RDDs zu Kanten- und Knoten-RDDs die als Tupel einen Graph beschreiben \cite{graphx}. Auf diesen abgeleiteten Datentypen sind übliche RDD- und Mengenoperationen wie \lstinline!subgraph!, \lstinline!diff!, u.a. für das Erstellen und Modifizieren von Graphen möglich.
Ein solcher Graph kann z.B. Webseitenrelationen, jeder Link ist eine Kante im Graph, jede Domain ein Knoten, darstellen.
Auf diesen Graphen sind Algorithmen wie PageRank oder das Auszählen von Dreiecken in Graphen ausführbar.

%%%%%%%%%%%%%%%%%%%%%
\section{MapReduce}
%%%%%%%%%%%%%%%%%%%%%

Spark basiert auf dem MapReduce-Paradigma, welches 2004 in einem Paper der Google-Mitarbeiter Dean und Ghemawat\cite{mapreduce2004,mapreduce2008} in Anlehnung an die aus funktionalen Programmiersprachen bekannten Map- und Reduce-Befehle eingeführt wurden.
Viele der benötigten Algorithmen wie Häufigkeitsanalyse oder Webseitengraphen hatten zuvor ihren eigenen Programmcode für die Kommunikation und Ausfallsicherheit im Cluster und folgten von der Struktur her einem simplem Schema: zuerst werden Eingabedaten datenparallel verarbeitet und anschließend werden die Zwischenergebnisse zu Endergebnissen reduziert.
Diese einfach gestrickten Algorithmen mussten aber trotzdem auf beträchtlichen Datenmengen auf Knoten hoher Ausfallwahrscheinlichkeit laufen. Diese Funktionalität, also die Kommunikation im Cluster und die Ausfallsicherheit. wurde so in eine MapReduce-Bibliothek ausgelagert.

MapReduce bezeichnet hierbei sowohl das Programmierparadigma als auch die Bibliothek, die diese ausfallsicher und parallelisiert zur Verfügung stellt.
Dafür definiert der Nutzer eine Map-Funktion die aus einer Liste jedes Schlüssel-Wert-Paar $(k,v)$ auf eine Liste aus neuen Schlüssel-Wert-Paaren abbildet:
\begin{equation}
    \mathbf{Map : }\left( k, v \right) \mapsto
    \left[ \left( l_1, x_1 \right), \ldots, \left( l_{r_k}, x_{r_k} \right) \right]
\end{equation}
$k$ und $l$ sind hierbei die Schlüssel und $v$ und $x$ die dazugehörigen Werte. Für $r_k=1$ erhält man den Grenzfall, dass jeder Schlüssel-Wert-Tupel auf exakt einen neuen Schlüssel-Wert-Tupel abgebildet wird.

Da es per Definition keine Abhängigkeiten zu anderen Schlüssel-Wert-Paaren für die Berechnung der Map-Funktion gibt, kann jedes Datum parallel berechnet werden.
Das MapReduce-Framework stellt außerdem Funktionen zum Einlesen von Daten zur Verfügung und verteilt diese automatisch an die Knoten des Clusters, wo diese dann parallel verarbeitet werden.

In einem impliziten Schuffle-Schritt werden alle Paare mit gleichem Schlüssel lokal gruppiert, d.h. es werden alle Paare z.B. mit Schlüssel $k_1=\mathrm{'Dresden'}$ auf dem gleichen Knoten im Cluster gesammelt.
Dieser Schritt kann also sehr kommunikationsaufwendig sein.
Der Nutzer kann nur über die Art der Schlüssel auf diesen Prozess Einfluss nehmen.
Man erhält als Ergebnis des Schuffle-Schritts einen Tupel aus einem Schlüssel und einer Liste an dazugehörigen Werten.
\begin{equation}
    \mathbf{Shuffle :}
    \left[ \left( k_1, x_1 \right), \ldots, \left( k_n, x_n \right) \right]
    \mapsto \left[ \left( l_1, \left[ y^1_1, \ldots, y^1_{r_1} \right] \right), \ldots,
    \left( l_m, \left[ y^m_1, \ldots, y^m_{r_m}) \right] \right) \right]
\end{equation}

Im nachfolgenden Reduce Schritt werden die Tupel aus Schlüssel und Werteliste mit einer vom Nutzer defineirten Reduce-Funktion abgebildet auf eine neue Werteliste.
Im einfachsten Fall kann die neue Werteliste nur ein Element enthalten, z.B. die Summe oder den Mittelwert der alten Werte.
\begin{equation}
    \mathbf{Reduce : }
    \left( l, \left[ y_1, \ldots, y_{s_l} \right] \right)
    \mapsto \left[ w_1, \ldots, w_{m_l} \right]
\end{equation}

Auch die Reduce-Operation kann also bezüglich einzigartiger Schlüssel datenparallel ausgeführt werden.
Anfangs war dieses Paradigma nur für einfache Beispiele wie Wortfrequenzanalysen gedacht, aber mittlerweile wurden auch komplexere Algorithmen wie z.B. Matrixmultiplikation\cite{matmulmapreduce} oder das Problem des Finden einer maximalen Überdeckung\cite{maxcovermapreduce} ins MapReduce-Schema portiert.\\

%%%%%%%%%%%%%%%%%%%%%
% MapReduce
%%%%%%%%%%%%%%%%%%%%%

Eine lange Zeit beliebte Implementation basierend auf dem MapReduce-Paper\cite{mapreduce2004} ist Apache Hadoop\cite{hadoop}.
Hadoop besteht aus einem verteilten Dateisystem, dem Hadoop Distributed File System (HDFS), und einer Bibliothek namens MapReduce, die Funktionen für das Rechnen auf den verteilten Daten anbietet.

Apache Spark\cite{spark} ist eine alternative Implementation des MapReduce-Modells.
Es versucht viele der Probleme von Hadoop zu beheben, bietet jedoch neben dem Zugriff vom lokalen Dateisystem auch den Zugriff von HDFS und weiteren verteilten Dateisystemen aus an.

Einer der Hauptvorteile von Spark ist der Geschwindigkeitsgewinn bei der iterativen Anwendung von Map und Reduce durch die Möglichkeit die Daten auch im Arbeitsspeicher, nicht nur auf der Festplatte, der Knoten zwischenzuspeichern.


%%%%%%%%%%%%%%%%%%%%%%%%%%%%%%%%%%%%%%%%%%%%%%%%%%%%%%%%%%%%%%%%%%%%%%%%%%%%%%%%
\section{Architektur}
%%%%%%%%%%%%%%%%%%%%%%%%%%%%%%%%%%%%%%%%%%%%%%%%%%%%%%%%%%%%%%%%%%%%%%%%%%%%%%%%

% Executor, Masternode
Für kleinere Tests kann Spark local auf einem Computer bzw. Knoten ausgeführt werden:
\begin{lstlisting}[numbers=none]
spark-shell --master local[*]
\end{lstlisting}\vspace{-1.5\baselineskip}
In diesem Beispielaufruf der interaktiven Spark-Shell werden so viele Threads genutzt wie es logische Kerne gibt.

Für das Ausführen auf einem Cluster ist jedoch das getrennte Starten von einem Spark Driver (Master) und mindestens einem Executor (Slave, Worker) notwendig, vgl. \autoref{fig:sparkcluster}.
Der Spark-Driver führt das geschriebene Spark-Programm aus und verteilt u.a. die für die Map-Funktion benötigten Daten an die Executoren, welche die zugewiesenen Berechnungen dann durchführen.

\begin{figure}
    \centerline{\includegraphics[width=0.8\linewidth]{cluster-overview.png}}
    \caption{Master/Slave-Struktur eines Spark-Clusters\cite{spark}}
    \label{fig:sparkcluster}
\end{figure}

Jeder Executor wird in einer eigenen Java Virtual Machine (JVM), also einem eigenen Prozess ausgeführt.
Normalerweise, und so auch hier, wird für jeden Knoten ein Executor-Prozess gestartet, welcher mit der in der \lstinline!SPARK_WORKER_CORES!-Umgebungsvariable definierten Anzahl an Threads arbeitet.
Für den hier vorgestellten Benchmark wird \lstinline!$SPARK_WORKER_CORES! identisch der Anzahl an Grafikkarten auf dem Knoten gewählt, sodass jeder Thread mit genau einer Grafikkarte arbeitet.\\

Ein RDD kann z.B. mit \lstinline!parallelize! erstellt werden \cite{sparkguide}.
\begin{lstlisting}[numbers=none,language=scala]
val distData = sc.parallelize( 1 to 10, 5 )
\end{lstlisting}\vspace{-1.5\baselineskip}
In diesem Beispiel wurden die Zahlen 1 bis 10 in ein RDD geladen.
Das zweite Argument gibt die Anzahl an Partitionen an, die das RDD nutzen soll.
Hier werden die Daten also in 5 Partitionen, in älteren Versionen von Spark auch Slices genannt, aufgeteilt.
Eine Partition ist wie ein Task zu verstehen, der von einem der Executoren ausgeführt wird.
Die Anzahl an Partitionen gibt damit eine Obergrenze für die Parallelisierbarkeit vor.
Um Gebrauch von einer optimalen Lastverteilung zu machen sollte man eine Größenordnung mehr Partitionen haben, als die Sparkinstanz logische Kerne bzw. Threads zur Verfügung hat.


%%%%%%%%%%%%%%%%%%%%%%%%%%%%%%%%%%%%%%%%%%%%%%%%%%%%%%%%%%%%%%%%%%%%%%%%%%%%%%%%
\section{Konfiguration von Spark auf einem Slurm-Cluster}
%%%%%%%%%%%%%%%%%%%%%%%%%%%%%%%%%%%%%%%%%%%%%%%%%%%%%%%%%%%%%%%%%%%%%%%%%%%%%%%%

Viele Cluster stellen schon einen Task-Scheduler für Multinutzerumgebungen, wie z.B. PBS (Portable Batch System) oder SLURM (Simple Linux Utility for Resource Management), zur Verfügung, um Rechenzeit auf dem Cluster möglichst effizient und gerecht zu verteilen.
Außerdem ermöglichen sie das verteilte Starten von Programmen, z.B. jene die mit MPI programmiert wurden.
Der für die Benchmarks genutzte Cluster, siehe \autoref{sct:taurus}, arbeitet mit SLURM.\\

Um Spark nutzen zu können, müssen zuerst Master- und Slave-Knoten gestartet werden.
Damit alle gleichzeitig gestartet werden, sollten sie innerhalb desselben Jobs initiiert werden.
Dies z.B. mit SLURMs \lstinline!--multi-prog! Option realisiert werden.
Diese erwartet als Argument einen Pfad zu einer Konfigurationsdatei, in der für jeden Rang ein möglicherweise anderes auszuführendes Programm angegeben ist.

Alternativ kann man auch anhand von der Umgebungsvariable \lstinline!SLURM_PROCID! im Skript entweder einen Master-Knoten oder einen Slave-Knoten starten.
Diese Möglichkeit wurde aufgrund der Übersichtlichkeit, weil damit alle Funktionalitäten in einem Skript vereint werden können, gewählt, siehe \lstinline!startSpark.sh!\attachfile{../startSpark.sh} und \lstinline!common/start_spark_slurm.sh!\attachfile{../common/start_spark_slurm.sh} \cite{scaromare}.

Auf dem Master-Knoten wird der Spark Driver mit
\begin{lstlisting}
"$SPARK_ROOT/bin/spark-class" org.apache.spark.deploy.master.Master \
    --ip $(hostname) --port 7077 --webui-port 8080 &
\end{lstlisting}\vspace{-1.5\baselineskip}
gestartet.
Alle anderen Knoten starten einen Executor-Prozess mit:
\begin{lstlisting}
"$SPARK_ROOT/bin/spark-class" org.apache.spark.deploy.worker.Worker \
    spark://$(scontrol show hostname $SLURM_NODELIST | head -n 1):7077
\end{lstlisting}\vspace{-1.5\baselineskip}
Hierbei wird vorausgesetzt, dass der Masterknoten, also jener für den \lstinline!$SLURM_PROCID=0! ist, der erste Knoten in \lstinline!$SLURM_NODELIST! ist. Dies wurde per assert vom Masterknoten aus auch geprüft und ist bei keinem der rund 50 Versuche fehlgeschlagen.

Wenn Spark gestartet ist, kann sich z.B. mit einer aktiven Eingabeaufforderung an den Master verbunden werden:
\begin{lstlisting}[language=bash]
export MASTER_ADDRESS=spark://$MASTER_IP:7077
spark-shell --master=$MASTER_ADDRESS
\end{lstlisting}\vspace{-1.5\baselineskip}
Die Umgebungsvariable \lstinline!MASTER_ADDRESS! wird automatisch vom \lstinline!startSpark.sh!-Skript \cite{scaromare} gesetzt.


%%%%%%%%%%%%%%%%%%%%%%%%%%%%%%%%%%%%%%%%%%%%%%%%%%%%%%%%%%%%%%%%%%%%%%%%%%%%%%%%
\chapter{Rootbeer}
\label{sct:rootbeer}
%%%%%%%%%%%%%%%%%%%%%%%%%%%%%%%%%%%%%%%%%%%%%%%%%%%%%%%%%%%%%%%%%%%%%%%%%%%%%%%%

Rootbeer\cite{pratt2012rootbeer} ist ein von Philip C. Pratt-Szeliga entwickeltes Programm welches das Schreiben von CUDA-Kerneln in Java ermöglicht.
Zum aktuellen Zeitpunkt, August 2016, hat Rootbeer leider noch Beta-Status und wurde seit ca. einem Jahr nicht weiterentwickelt\cite{rootbeergithub}.

Der Rootbeerquellcode enthält Gerüste, um den Rootbeerkernel nicht nur nach CUDA, sondern auch nach OpenCL zu übersetzen, aber diese Funktionalität ist noch nicht fertiggestellt.

Andere Ansätze für Grafikkartenprogrammierung innerhalb von Java bzw. Scala stellen simple, teilweise direkt von den Headern der OpenCL-Spezifikation kompilierte, oder aber auch komplexere objektorientierte Java-Schnittstellen für OpenCL und CUDA zur Verfügung, so z.B. JogAmp JOCL\cite{jogampcl}, JOCL\cite{jocl}, JavaCL\cite{javacl} und jCUDA\cite{jcuda} bzw. ScalaCL\cite{scalacl} und Firepile\cite{firepile} für Scala.
Diese erfordern jedoch die Kenntnis von OpenCL bzw. CUDA und der Nutzer muss sich selbstständig um Serialisierung und Host-GPU-Transfers kümmern.
ScalaCL und Firepile sind außerdem auch noch in der Entwicklung.

Desweiteren gibt es fertige Bibliotheken, die auf Grafikkarten portierte Funktionen zur Verfügung stellen, welche jedoch nur eine beschränkte Einsetzbarkeit haben.

Eine Alternative zu Rootbeer stellt das ursprünglich von AMD entwickelte Aparapi\cite{aparapi}, welches seit 2011 Open Source ist, dar.
Wie in Rootbeer ist es möglich in Java Kernel zu schreiben, die von Aparapi in OpenCL übersetzt werden.
Jedoch werden nur eindimensionale Felder unterstützt, im Gegensatz zu Rootbeer welches auch komplexe Objekte serialisieren kann.\\


%%%%%%%%%%%%%%%%%%%%%%%%%%%%%%%%%%%%%%%%%%%%%%%%%%%%%%%%%%%%%%%%%%%%%%%%%%%%%%%%
%\section{Funktionsweise}
%%%%%%%%%%%%%%%%%%%%%%%%%%%%%%%%%%%%%%%%%%%%%%%%%%%%%%%%%%%%%%%%%%%%%%%%%%%%%%%%


%%%%%%%%%%%%%%%%%%%%%%%%%%%%%%%%%%%%%%%%%%%%%%%%%%%%%%%%%%%%%%%%%%%%%%%%%%%%%%%%
%\subsection{Nutzung}
%%%%%%%%%%%%%%%%%%%%%%%%%%%%%%%%%%%%%%%%%%%%%%%%%%%%%%%%%%%%%%%%%%%%%%%%%%%%%%%%

Für die Nutzung von Rootbeer ist das Java SE Development Kit 7 und das NVIDIA CUDA Toolkit notwendig.
Rootbeer funktioniert noch nicht mit JDK 8 und unterstützt offiziell nur JDK 6 \cite{rootbeerjdk6}, aber in den hier durchgeführten Beispielen gab es keine Probleme mit JDK 7.

Zuerst muss der Nutzer das \lstinline!org.trifort.rootbeer.runtime.Kernel!-Interface implementieren und die kompilierten Klassen als JAR-Archive nochmals an den Rootbeer-Compiler geben, siehe \autoref{sct:kernelimplementation} und \autoref{sct:compilation}.
Es ist zu beachten, dass die an Rootbeer übergebene Datei auf \texttt{.jar} enden muss, insbesondere führen Dateinamen wie \texttt{gpu.jar.tmp} zu einer Fehlermeldung.\\

Der Rootbeer-Compiler nutzt dann Soot\cite{sootsite,sootretrospective}, um den Java Bytecode der Kernel-Implementation in Jimple zu übersetzen.
Jimple ist eine vereinfachte Zwischendarstellung von Java-Bytecode in Drei-Address-Code mit nur 15 verschiedenen Befehlen. Java-Bytecode hingegen hat ca. 200 verschiedene Befehle.
Der Jimple-Code wird dann analysiert und in CUDA-Quellcode übersetzt, welcher dann mit dem NVIDIA-Compiler kompiliert wird.
All das geschieht automatisch, aber die Zwischenschritte kann man zur Fehlersuche unter Linux in \lstinline!$HOME/.rootbeer/! einsehen.
Die so erstellte cubin-Datei wird zusammen mit \texttt{Rootbeer.jar} dem JAR-Archiv des Nutzers hinzugefügt.

Die zweite große Vereinfachung, die Rootbeer zur Verfügung stellt, ist die Automatisierung des Datentransfers zwischen GPU und CPU.
Das besondere hierbei ist, dass Rootbeer die Nutzung von beliebigen, also insbesondere auch nicht-primitiven Datentypen erlaubt.
Diese Datentypen serialisiert Rootbeer automatisch und unter Nutzung aller CPU-Kerne und transferiert sie danach auf die Grafikkarte.

Diese zwei Vereinfachungen obig machen die erste Nutzung von Rootbeer verglichen zu anderen Lösungen sehr einfach, sodass Rootbeer insbesondere für das Erstellen von Prototypen günstig ist.
In Kontrast dazu ist es jedoch auch möglich sehr nah an der Grafikkarte zu programmieren.
Dafür kann man mit Rootbeer manuell die Kernel-Konfiguration angeben, mehrere GPUs ansprechen, shared memory nutzen und auch über die \lstinline!RootbeerGpu!-Klasse CUDA-Befehle wie \lstinline!syncthreads! benutzen.

\chapter{Eigene Implementierung}
\label{sct:implementation}

\section{Kompilierung}

%Versucht man mit \lstinline!Rootbeer.jar! eine jar-Datei, welche auch Scala-Klassen beinhaltet, in GPU-Code umsetzen zu lassen, z.B. auf diese Weise:
%
%\begin{lstlisting}[language=bash]
%scalac $< -classpath $(ROOTBEER_ROOT)/Rootbeer.jar:. -deprecation
%javac $< -classpath $(ROOTBEER_ROOT)/Rootbeer.jar:.
%scalac $< -classpath $(ROOTBEER_ROOT)/Rootbeer.jar:. -deprecation
%jar -cvfm MontePiGPU.jar.tmp.jar manifest.txt *.class
%java -jar Rootbeer.jar MontePiGPU.jar.tmp.jar MontePiGPU.jar -64bit -computecapability=sm_30
%\end{lstlisting}
%
%Dann kommt es bei soot zu einer Runtime Exception:
%
%\begin{lstlisting}
%java.lang.RuntimeException: cannot get resident body for phantom class : <MonteCarloPiKernel$$anonfun$gpuMethod$1: void <init>(MonteCarloPiKernel,scala.runtime.ObjectRef)>; maybe you want to call c.setApplicationClass() on this class!
%	at soot.SootMethod.retrieveActiveBody(SootMethod.java:316)
%	at soot.rbclassload.RootbeerClassLoader.loadScene(RootbeerClassLoader.java:857)
%	at soot.rbclassload.RootbeerClassLoader.loadNecessaryClasses(RootbeerClassLoader.java:320)
%	at org.trifort.rootbeer.entry.RootbeerCompiler.setupSoot(RootbeerCompiler.java:198)
%	at org.trifort.rootbeer.entry.RootbeerCompiler.compile(RootbeerCompiler.java:219)
%	at org.trifort.rootbeer.entry.RootbeerCompiler.compile(RootbeerCompiler.java:213)
%	at org.trifort.rootbeer.entry.Main.run(Main.java:208)
%	at org.trifort.rootbeer.entry.Main.main(Main.java:244)
%\end{lstlisting}
%
%"Phantom classes are implicitly created models of classes that do not exist on Soot's classpath. To get rid of this problem, make sure that the class you are referencing is on Soot's classpath." \url{http://stackoverflow.com/questions/27383832/why-is-soot-always-saying-that-the-class-i-want-to-load-is-a-phantom-class-even}
%
%Nachdem nun also \lstinline!spark-library.jar! hinzugefügt wurde, verschwindet diese exception, aber Rootbeer findet die Kernel-Implementierung nicht:
%\begin{lstlisting}
%There are no kernel classes. Please implement the following interface to use rootbeer:
%org.trifort.runtime.Kernel
%\end{lstlisting}
%
%Dies könnte daran liegen, das \lstinline!MonteCarloPiKernel.scala! scala benutzt. Dadurch wird für \texttt{soot} verschleiert, dass die Klasse \texttt{Kernel} implementiert.
%Eine Implementierung in java löst das Problem jedoch nicht. Vermutlich weil \texttt{soot} oder \texttt{Rootbeer} nur Klassen analysieren die von der main-Klasse aufgerufen werden. Da die  main-Klasse aber in scala geschrieben ist, ist der Aufruf womöglich für \texttt{soot} nicht sichtbar, sodass \texttt{soot} \texttt{MonteCarloPiKernel.class} als unbenutzt einstuft und ignoriert. Ein Hinweis darauf liefert diese Ausgabe von Soot bei der Ausführung von Rootbeer:
%\begin{lstlisting}
%Total loaded classes: 557
%Total loaded methods: 0
%\end{lstlisting}
%Bei der Java-Version funktioniert die Kompilierung mit Rootbeer und die Ausgabe ist:
%\begin{lstlisting}
%Total loaded classes: 1428
%Total loaded methods: 4385
%\end{lstlisting}
%Eine mögliche Lösung wäre es eine Dummy-Java-Klasse zu schreiben die alle Kernel-Implementierungen aufruft und als main-class im Manifest festgelegt ist.
%Das Problem hierbei ist jedoch, dass Rootbeer Serialisierungs- und Deserialisierungscode in die aufrufende Klasse injeziert. Damit dieser injezierte Code also nicht in die Dummy-Klasse eingesetzt wird, braucht es noch eine Wrapper-Klasse, welche die Kernel-Implementierung aufruft und von der Dummy-Klasse aufgerufen wird. Dies ist \texttt{MonteCarloPi.java}.
%
%Nachdem die Kompilierung mit \texttt{Rootbeer} dann erfolgt ist, muss die Manifest-Datei \texttt{META-INF/MANIFEST.MF} dann nur noch abgeändert werden, damit sie nicht mehr auf die Dummy-Klasse verweist. Durchgeführt wurde es ähnlich, jedoch wurden alle nicht für Rootbeer relevanten Klassen erst nach dem Rootbeer-Aufruf zusammen mit der neuen Manifest-Datei in eine jar-Datei zusammengeführt, siehe Abb.\ref{fig:compilation}.


\begin{figure}[H]
	\centering
	\begin{minipage}{\linewidth}
		\includegraphics[width=\linewidth]{compile-structure-deu.pdf}
	\end{minipage}
	\caption{Kompilationsschema mit Kommadozeilenbefehlen und Zwischenstati.}
	\label{fig:compilation}
\end{figure}

!!! Problem: Hab mehrere Fehler gefunden deren Kenntnis möglicherweise eine Vereinfachung des Schemas bedeutet. Da bin ich noch am rumspielen, daher ist das halbfertig.

\section{Probleme}
Implementierung:
    Was ist bei GPUs zu beachten ( Seeds, 64-Bit )
    Was ist bei Rootbeer zu beachten?
      - private Variablen werden wirklich immer per memcpy hin und her transportiert.
      - muss nicht auf ungerade Kernel-Zahl achten, werden automatisch aussortiert

%%%%%%%%%%%%%%%%%%%%%%%%%%%%%%%%%%%%%%%%%%%%%%%%%%%%%%%%%%%%%%%%%%%%%%%%%%%%%%%%
\chapter{Leistungsanalyse}
\label{sct:benchmarks}
%%%%%%%%%%%%%%%%%%%%%%%%%%%%%%%%%%%%%%%%%%%%%%%%%%%%%%%%%%%%%%%%%%%%%%%%%%%%%%%%


%\subsection{Pseudozufallszahlgeneratoren}
%\section{Vergleich - Alternativen}


%%%%%%%%%%%%%%%%%%%%%%%%%%%%%%%%%%%%%%%%%%%%%%%%%%%%%%%%%%%%%%%%%%%%%%%%%%%%%%%%
\section{Testsysteme}
%%%%%%%%%%%%%%%%%%%%%%%%%%%%%%%%%%%%%%%%%%%%%%%%%%%%%%%%%%%%%%%%%%%%%%%%%%%%%%%%


%%%%%%%%%%%%%%%%%%%%%%%%%%%%%%%%%%%%%%%%%%%%%%%%%%%%%%%%%%%%%%%%%%%%%%%%%%%%%%%%
\subsection{System 1: Heimsystem}
\label{sct:system1}
%%%%%%%%%%%%%%%%%%%%%%%%%%%%%%%%%%%%%%%%%%%%%%%%%%%%%%%%%%%%%%%%%%%%%%%%%%%%%%%%


Aufgrund der leichten Verfügbarkeit und als Beispiel für Grafikkartenbeschleuniger im nicht-professionellen Verbrauchersegment, wurden einige Tests auf einem herkömmlichen Arbeitsplatzrechner ausgeführt:

\begin{table}[H]
	\begin{center}\begin{tabularx}{0.9\linewidth}{r|X}
	Prozessor & Intel(R) Core(TM) i3-3220 (Ivy-Bridge), \SI{3.30}{\giga\hertz}, 2 Kerne (4 durch SMT), AVX\cite{ark3220} \\
	Arbeitsspeicher & $4\times \SI{4}{\gibi\byte}$ DDR3 \SI{1600}{\mega\hertz} CL9\cite{corsair4gbddr3}\\
	\hline
	Grafikkarte & GigaByte GTX 760 WindForce 3X Overclocked, Codename: GK104-225-A2 (Kepler), 1152 CUDA-Kerne, \SI{1085}{\mega\hertz} ( \SI{1150}{\mega\hertz} Boost), \SI{2}{\gibi\byte} GDDR5-VRAM mit \SI{1502}{\mega\hertz} PCIe-3.0-x16-Schnittstelle\cite{gigabytegtx760,gtx760,nvidiakepler}
	\end{tabularx}\end{center}
	\caption{Heimsystemkonfiguration}
\end{table}

Die Maximalleistung in der Berechnung von Fließkommazahlen einfacher Genauigkeit (SPFLO) im Verhältnis einer Multiplikation zu einer Addition beträgt
\begin{align}
	\SI{3.30}{\giga\hertz} \cdot
	2\,\text{Kerne} \left(
		1\,\frac{ \text{AVX ADD Einheit} }{ \text{Kern} } +
		1\,\frac{ \text{AVX MUL Einheit} }{ \text{Kern} }
	\right) \cdot
	8\,\frac{ \text{SPFLO} }{ \text{AVX Einheit} } \\
	= 105.6\,\text{GSPFLOPS}
\end{align}
für den Prozessor. Informationen zur Architektur, wie die Anzahl an AVX-Einheiten wurde aus Ref.\cite{cesga} entnommen.

Die Maximalleistung der Grafikkarte beträgt:
\begin{align}
	\SI{1.085}{\giga\hertz} \cdot 1152 \text{CUDA-Kerne} \cdot
	1 \frac{ \text{FMA-Einheit} }{ \text{CUDA-Kern} } \cdot
	2 \frac{ \text{SPFLO} }{ \text{FMA-Einheit} }
	= 2500\,\text{GSPFLOPS}
\end{align}
Zugegeben, es wurde beim Prozessor gespart und bei der Grafikkarte nicht, aber der Geschwindigkeitsunterschied von 24x begründet dennoch das Interesse daran Grafikkarten zu nutzen, auch wenn es bei Grafikkarten mehr zu beachten gibt, um diese Maximalleistung erhalten zu können.



%%%%%%%%%%%%%%%%%%%%%%%%%%%%%%%%%%%%%%%%%%%%%%%%%%%%%%%%%%%%%%%%%%%%%%%%%%%%%%%%
\subsection{System 2: Taurus}
\label{sct:taurus}
%%%%%%%%%%%%%%%%%%%%%%%%%%%%%%%%%%%%%%%%%%%%%%%%%%%%%%%%%%%%%%%%%%%%%%%%%%%%%%%%


Für Skalierungstests wurde einer der Hochleistungsrechner der TU-Dresden, ein Bull HPC-Cluster mit dem Namen Taurus, benutzt. Der Bau der ersten Phase von Taurus war 2013 abgeschlossen\cite{taurusnutzerschulung}. Zum Zeitpunkt der Nutzung (2015/2016) waren alle Knoten von Phase 1 schon in die 2015 fertiggestellt\cite{heisehrsk2} Phase 2 integriert wurden\cite{doctudtaurushardware} und werden nun beide unter dem Namen Taurus zusammengefasst.

\begin{figure}[H]
	\centering
	\begin{minipage}{0.5\linewidth}
		\includegraphics[width=\linewidth]{taurusphase1}
	\end{minipage}\begin{minipage}{0.5\linewidth}
		\includegraphics[width=\linewidth]{taurusphase1island1interconnects}
	\end{minipage}
	\caption{\textbf{Links:} Übersicht Taurus Phase 1. \textbf{Rechts:} Schema der Topologie von Insel 2, auf der ausschließlich gerechnet wurde. Die Bilder wurden übernommen aus Ref.\cite{taurusnutzerschulung}}
	\label{fig:taurusphase1}
\end{figure}

Gerechnet wurde auf Insel 2 von Taurus, vgl. Tabelle~\ref{tbl:island2}. Wenn nicht anders erwähnt, dann beziehen sich Benchmarks auf die Tesla K20x Knoten.

\begin{table}[H]
	\begin{tabularx}{0.9\linewidth}{r|Y|Y}
		& \textbf{Phase 1} & \textbf{Phase 2} \\
		\hline
		Knoten          & 44 & 64 \\
		Hostnamen       & taurusi2[001-044] & taurusi2[045-108] \\
		Prozessor       & 2x Intel Xeon CPU E5-2450 (8 Kerne) @ 2.10GHz, MultiThreading deaktiviert, AVX, 2x 268.8\,GSPFLOPS
						& 2x Intel(R) Xeon(R) CPU E5-2680 v3 (12 Kerne) @ 2.50GHz, MultiThreading deaktiviert, AVX2 insbesondere FMA3\cite{ark2680v3}, 2x 537.6\,GSPFLOPS \\
		GPU 			& 2x NVIDIA Tesla K20x & 4x NVIDIA Tesla K80 \\
		Arbeitsspeicher & \SI{32}{\gibi\byte} & \SI{64}{\gibi\byte} \\
		Festplatte      & \SI{128}{\gibi\byte} SSD & \SI{128}{\gibi\byte} SSD
	\end{tabularx}
	\caption{Zusammensetzung Insel 2 von Taurus\cite{doctudtaurussystem}}
	\label{tbl:island2}
\end{table}

\begin{table}[H]
	\begin{tabular}{r|c|c}
	& \textbf{K20x} & \textbf{K80} \\
	\hline
	Chip       & GK110 & GK210 \\
	Takt       & \SI{0.732}{\giga\hertz} & \SI{0.560}{\giga\hertz} \\
	CUDA-Kerne & 2688 & 4992 \\
	Speicher   & \SI{6}{\gibi\byte}, GDDR5 384 Bit Busbreite
	           & $2\times\SI{12}{\gibi\byte}$, GDDR5 384 Bit Busbreite \\
    Bandbreite & \SI{250}{\giga\byte\per\second}
	           & $2\times\SI{240}{\giga\byte\per\second}$              \\
%	\begin{minipage}{2.5cm}\begin{flushright}Theoretische\\Spitzenleistung\end{flushright}\end{minipage} & 3935\,GSPFLOPS, 1312\,GDPFLOPS
%					& 5591\,GSPFLOPS, 1864\,GDPFLOPS \\
	Theoretische    & 3935\,GSPFLOPS & 5591\,GSPFLOPS \\
	Spitzenleistung & 1312\,GDPFLOPS & 1864\,GDPFLOPS
	\end{tabular}
	\caption{Spezifikationen der Kepler-Grafikkarten von Taurus\cite{nvidiakepler,k20anandtech}}
	\label{tbl:k20k80}
\end{table}

Zu den Spitzenleistungen in Tabelle~\ref{tbl:k20k80} sei angemerkt, dass jeder CUDA-Kern einfache Fließkommagenaugigkeit berechnet und auf drei CUDA-Kerne eine Doppelpräzisionseinheit kommt, wodurch sich die DFLOPS berechnen.
Die K80 hat außerdem einen Boost-Modus mit \SI{0.875}{\giga\hertz}, also einer Leistungssteigerung von $1.56$.

\section{Monte-Carlo-Simulation verschiedener Implementationen}

In Abb.\ref{fig:montepiworkloadscaling} wurde die Ausführungszeit von Monte-Carlo-Simulationen in verschiedenen Programmiersprachen über die Anzahl an Monte-Carlo-Iterationen gemessen. Gemessen wurde die Zeit mit dem Linux \texttt{time}-Befehl und zwar die \texttt{real}-Zeit. Es fällt auf, dass alle Versionen eine Initialisierungszeit haben. Bei der C++-Version beträgt diese jedoch nur knapp \SI{30}{\milli\second}, während die reine Java-Version schon ca. \SI{70}{\milli\second} benötigt. Die Nutzung von Scala erhöht dies schon auf \SI{200}{\milli\second} und die Nutzung von Rootbeer führt eine weitere Initialisierungszeit von \SI{430}{\milli\second} ein, sodass für wenig Iterationen die Rootberversion bis zu 20x langsamer sind. Erst für 10 Milliarden Iterationen beginnt die Initialisierungszeit im Vergleich zur Rechenzeit vernachlässigbar zu werden, sodass aber da die Lastenskalierung ein lineares Verhalten annimmt. Bei der C++-Version ist dies schon bei ca. 100 Millionen Iterationen der Fall.
\begin{figure}[H]
	\centering
	\begin{minipage}{0.5\linewidth}
		\includegraphics[width=\linewidth]{benchmarks-workload-scaling}
	\end{minipage}\begin{minipage}{0.5\linewidth}
		\includegraphics[width=\linewidth]{benchmarks-workload-scaling-taurus2}
	\end{minipage}
	\caption{Benötigte Ausführungszeit der Monte-Carlo Pi-Berechnung in Abhängigkeit von der Anzahl an Iterationen. Getestet auf \textbf{links:} System 1 und \textbf{rechts:} System 2 (Taurus), siehe Kapitel~\ref{sct:system1}}
	\label{fig:montepiworkloadscaling}
\end{figure}
Im Programmausdruck~\ref{lst:montemainloop} ist die Hauptschleife, die hauptsächlich die Arbeitslast generiert, zu sehen. Es handelt sich also für jede der zwei Zufallszahlen um eine Multiplikations und zwei Divisionen und dann nochmals zwei Multiplikationen für die Berechnung des Quadrat des Radius, also zusammen acht Operationen pro Iteration und neun Operationen für Iterationen die im Kreis liegen. Dies tritt für $\frac{\pi}{4}=0.785\%$ der Fälle auf, das heißt die Rechenlast, definiert als die Anzahl an arithmetischen Operationen $N_\text{Op}$ sollte sich wie folgt aus der Anzahl an Iterationen $N$ berechnen:
\begin{align}
	N_\text{Op}
	= \left[ \frac{\pi}{4}\cdot 9 + \left( 1-\frac{\pi}{4} \right)\cdot 8 \right] N
	= 8.8\cdot N
\end{align}

Zuletzt ist aus dem Plot abzulesen, dass für große Lasten wie z.B. für drei Milliarden Iterationen die Versionen, die von Grafikarten Gebrauch machen, um einen Faktor $140$ (Scala) bis $320$ (C++) schneller sind als die Java Version, die auf einem Prozessor-Kern ausgeführt wird. Dieser große Unterschied zwischen CPU und GPU lässt darauf schließen, dass Java weder AVX noch mehrere Kerne gleichzeitig nutzt. Das heißt für die CPU-Versionen wäre für das Testsystem 1 noch ein Geschwindigkeitsgewinn von $8\,\frac{ \text{Op} }{ \text{AVX-Einheit} } \cdot 2\,\text{Kerne}$ erreichbar. Dies würde den Geschwindigkeitsunterschied von $320$ auf $20$ reduzieren, was in Übereinstimmung mit dem Verhältnis der Peakflops aus Kapitel~\ref{sct:system1} wäre.

\begin{lstlisting}[language=Java,caption={Hauptschleife der Monte-Carlo Pi-Berechnung},label=lst:montemainloop]
for ( int i = 0; i < dnDiceRolls; ++i )
{
    dRandomSeed = (int)( (randMagic*dRandomSeed) % randMax );
    float x = (float) dRandomSeed / randMax;
    dRandomSeed = (int)( (randMagic*dRandomSeed) % randMax );
    float y = (float) dRandomSeed / randMax;
    if ( x*x + y*y < 1.0 )
        nHits += 1;
}
\end{lstlisting}
Dieses Branching verändert also nicht das Skalierverhalten linear mit $N$, sondern führt nur zu einem veränderten Faktor. Daher ist in Abb.\ref{fig:montepiworkloadscaling} lineares Verhalten zu beobachten.



%%%%%%%%%%%%%%%%%%%%%%%%%%%%%%%%%%%%%%%%%%%%%%%%%%%%%%%%%%%%%%%%%%%%%%%%%%%%%%%%
\section{Monte-Carlo-Simulation mit Spark + Rootbeer}
%%%%%%%%%%%%%%%%%%%%%%%%%%%%%%%%%%%%%%%%%%%%%%%%%%%%%%%%%%%%%%%%%%%%%%%%%%%%%%%%


In Abbildung~\ref{fig:montepistrongscaling} ist die Laufzeit über die Anzahl an Kernen bzw. Grafikkarten dargestellt. Aus gründen der Rechenzeit wurde für eine Anzahl von $N=1,2,4,8,16,24,32$ Knoten die Leistungsanalyse für $4(N-1)$ bis $4N$ Kerne/Grafikkarten durchgeführt. Dies ist vor allem in der Abbildung rechts zu sehen, wo die Messpunkte immer in fünfer-Gruppen auftreten.

Zwar nicht in der Abbildung dargestellt, wurde auch für $4N+1$ und $4N+2$ gemessen. Das heißt Spark hat zwei mehr Slices zu verarbeiten als es Kerne gibt. Dies führt dazu, dass zwei Kerne doppelt so viel Arbeit wie der Rest haben, wodurch es zu einem sprunghaften Anstieg der Laufzeit kommt. Aus diesem Grund ist es normalerweise besser, wenn die Anzahl an Slices viel größer als die Anzahl an Kernen bzw. Grafikkarten wäre. Dies würde aber die ohnehin schon recht hohe Mindestproblemgröße noch einmal um ein, zwei weitere Größenordnungen erhöhen.

\begin{figure}[H]
	\centering
	\begin{minipage}{0.5\linewidth}
		\includegraphics[width=\linewidth]{cluster-strong-scaling-gpu}
	\end{minipage}\begin{minipage}{0.5\linewidth}
		\includegraphics[width=\linewidth]{cluster-strong-scaling-cpu}
	\end{minipage}
	\caption{Benötigte Ausführungszeit der Monte-Carlo Pi-Berechnung in Abhängigkeit von der Anzahl an \textbf{links:} Kernen und \textbf{rechts:} Grafikkarten. Getestet auf System 2, siehe Kapitel~\ref{sct:taurus}}
	\label{fig:montepistrongscaling}
\end{figure}

Weiterhin auffällig sind zufällig auftretende Spitzen in sowohl im CPU als auch im GPU-Benchmark. Z.b. für eine Arbeitslast von $2^33$ Iterationen pro Slice auf 24 Knoten also 92 bis 96 Grafikkarten schwankt die Ausführungszeit zwischen \SI{30}{\second}, \SI{45}{\second} und \SI{68}{\second}. Möglicherweise liegt dies an einer ungünstigen Verteilugn der Slices auf die Knoten. Die vorliegende Version nimmt eine lineare Verteilung an, sodass Slice 0 auf Grafikkarte 0 von Knoten 0 rechnet, während Slice 1 auf Grafikkarte 1 von Knoten 0 und Slice 4 auf Grafikkarte 0 von Knoten 1 rechnet. Die Zeit die eine Grafikkarte für diese Last benötigt ist \SI{22}{\second}. Es ist also wahrscheinlich, dass zumindest im Fall der \SI{68}{\second} drei verschiedene Prozesse auf einem Knoten diesselbe Grafikkarte anfordern. Diese Vermutung wurde getestet, indem jeder Slice seinen Hostnamen ausgeben soll. Mit dem Skript aus Listing~\ref{lst:start_spark_slurm.sh} ist dies schnell interaktiv getestet:
\begin{lstlisting}[language=scala]
startSpark --time=04:00:00 --nodes=5 --partition=west --gres= --cpus-per-task=12
spark-shell --master=$MASTER_ADDRESS
scala> import java.net.InetAddress
scala> sc.parallelize( 1 to 5*12, 5*12 ).map( (x) => { Thread.sleep(10); x+" : "+InetAddress.getLocalHost().getHostName() } ).collect().foreach( println )
\end{lstlisting}\vspace{-1.5\baselineskip}
Es ist also wie vermutet: die Verteilung ist nicht linear, sondern eher verzahnt, aber eigentlich zufällig.

Eine Lösung dieses Problems ist schwierig, da die Verteilung der Slices von Spark auf die Worker-Knoten opak erfolgt und es auch schwierig ist mit CUDA und vor allem mit Rootbeer herauszufinden welche der verfügbaren Grafikkarten in Benutzung ist. Ein ändern des Compute-Modus in einen Thread- oder Prozess-exklusiven Modus mittels
\begin{lstlisting}[language=bash]
nvidia-smi --compute-mode=EXCLUSIVE_PROCESS
\end{lstlisting}
ist auf Grund fehlender Berechtigungen im Cluster nicht möglich. In diesem Modus würde der Versuch eine schon in Benutzung seiende Grafikkarte anzusprechen in einer ''GPU device not available''-Fehlermeldung enden. Womöglich ist es gar nicht möglich dies über Rootbeer aus abzufangen.

Die Spitzen im CPU-Benchmark lassen sich dadurch jedoch nicht erklären, da sie für sehr Hohe Arbeitslasten kleiner verschwindet klein werden. Es handelt sich also wahrscheinlich eher um zufällige Initialisierungsoffsets oder Kommunikationslatenzen. Sie sind ungefähr \SI{3}{\second} groß, womit Kommunikationslatenz sehr unwahrscheinlich sind, da \lstinline!ping taurusi2063! als Beispiel Latenzen im Bereich von \SI{200}{\micro\second} misst. Es sei hier angemerkt, dass es in den 343 Testläufen für jeweils CPU und GPU zu zwei Fällen kam, in denen ein Job über das Spark-Web-Interface manuell beendet werden musste, da sie schon mehrere Minuten ohne Fortschritt liefen. Möglicherweise war dies aber auch ein Sympton der GPU-Konflikte pro Knoten.

In Abbildung~\ref{fig:montepistrongscaling} ist für große Arbeitslasten wie zu erwarten ein nahezu konstantes Verhalten über erhöhte Parallelisierung abzulesen. Für kleine Arbeitslasten ist eine schwache monotone Abhängigkeit zu beobachten. Möglicherweise ist dies die Zeit, die eine Reduktion über 32 Knoten länger braucht als z.B. über vier Knoten.

Die Benchmarks wurden für eine Grafikkarte pro Knoten wiederholt. Außerdem wurde leider festgestellt, dass die Benchmarks mit nur 384 Threads pro Grafikkarte ausgeführt wurden. Dies wurde geändert auf eine automatische Bestimmung, die zu ungefähr 20000 Threads führen sollte, womit Pipelining genutzt werden kann, sodass ein Faktor von 30 und mehr an Speedup zu erwarten ist.

\chapter{Zusammenfassung}

 - Ausblick(Nutzbarkeit, Anwendungsfälle, Deep Learning)
 - deep learning - Ähnlichkeit (assembling)

 - mehr cores



%\setlength{\headheight}{27pt} % needed because zih-template is buggy as *** and the chapter title is too long for the standard header
%\setlength{\headheight}{17pt}

\appendix

\chapter{Standardabweichung des Mittelwertes}
\label{apx:meanerror}

Dieses Kapitel richtet sich leicht nach Ref.\cite{meanerrorulm}. Sei $f\equiv(f_i)$ eine Folge von $N$ Stichproben aus einer Zufallsverteilung mit einem Mittelwert $\mu$ und $\mu_N$ der empirische Mittelwert dieser Folge. Die Differenz $\Delta_N:=\mu-\mu_N$ wird nun abgeschätzt mit der Standardabweichung einer Folge von empirischen Mittelwerten $\left(\mu_{N,k}\right)$ die alle mit (sehr wahrscheinlich) verschiedenen Folgen bzw. Vektoren $(f_i)$ gebildet seien.

Sei nun $\langle \cdot \rangle_N: \mathbb{S}\times \mathbb{R}^N \rightarrow \mathbb{R}^N, \langle (f_i)_k \rangle_N:= \frac{1}{N}\sum\limits_{i=0}^N {f_i}_k$ der empirische Mittelwert und $E(\cdot):\mathbb{S}\rightarrow \mathbb{R}, E\left( x_k \right):= \lim\limits_{N\rightarrow \infty} \frac{1}{N} \sum\limits_{k=0}^N f_k$ der Erwartungswert über eine unendlich Folge aus einem beliebigen statistischen Werten für begrenzte Folgen $(f_i)$. Hierbei ist $\mathbb{S}$ der Raum der Folgen. Aus der Definition der beiden Mittelwerte wird klar, dass man die Summen und damit die Bildung der Mittelwerte vertauschen kann, sofern ein Grenzwert existiert. Dies wird in Gleichung~\ref{eq:EabToEaEb-pre} angewandt.
%\\  = \frac{1}{N} E\left( \sigma_k \right)
%    + E\left( \left( \mu_{N,k}-\mu \right)
%      \frac{1}{N} \sum\limits_{j=0,j\neq i}^N \left( f_{jk}-\mu \right) \right)
%\\  \approx \frac{1}{N} E\left( \sigma_k \right) +
%      E\left( \frac{1}{N} \sum\limits_{i=0}^N \left( f_{ik}-\mu \right)
%      \left( \mu_{N,k'} - \mu \right) \right) \\
%	= \frac{\sigma}{N} +  E\left( \left(\mu_{N,k} - \mu \right)
%	  \left( \mu_{N,k} - \mu \right) \right)
%   = \frac{\sigma}{N} +  E\left( \mu_{N,k} - \mu \right)
%	  E\left( \mu_{N,k} - \mu  \right) \\
%	= \frac{\sigma}{N} +  \left( \mu - \mu \right)\left( \mu - \mu  \right)
\begin{align}
	\sigma_{\mu_N}^2
   := E\left( \left( \mu_N - \mu \right)^2 \right)
   := E\left( \left( \langle f_{ik} \rangle_N - \mu \right)^2 \right)
	= E\left( \langle f_{ik} - \mu \rangle_N^2 \right)
\\  = E\left(
	  \left( \frac{1}{N} \sum\limits_{i=0}^N \left(f_{ik}-\mu\right) \right)
	  \left( \frac{1}{N} \sum\limits_{i=0}^N \left(f_{ik}-\mu\right) \right)
	  \right)
	= E\left( \frac{1}{N^2} \sum\limits_{i=0}^N \sum\limits_{j=0}^N
		      \left( f_{ik}-\mu \right) \left( f_{jk}-\mu \right) \right)
    \\
	\label{eq:EabToEaEb-pre}
	= \frac{1}{N} E\left(
		 \frac{1}{N} \sum\limits_{i=0}^N \left( f_{ik}-\mu \right)^2
      \right)
    + E\left( \frac{1}{N} \sum\limits_{i=0}^N \left( f_{ik}-\mu \right)
      \frac{1}{N} \sum\limits_{j=0,j\neq i}^N \left( f_{jk}-\mu \right) \right)
    \\
	\label{eq:EabToEaEb-post}
    = \frac{1}{N} E\left( \sigma_k \right)
    + \frac{1}{N} \sum\limits_{i=0}^N
      \frac{1}{N} \sum\limits_{j=0,j\neq i}^N
      \left( \underbrace{E\left(f_{ik}\right)}_{=\mu}-\mu \right)
      \left( \underbrace{E\left(f_{jk}\right)}_{=\mu}-\mu \right)
	= \uuline{ \frac{\sigma}{N} }
\end{align}
Man beachte, dass der Schritt in Gl.\ref{eq:EabToEaEb-pre}-\ref{eq:EabToEaEb-post} nur möglich ist, wenn $f_i$ unabhängig von $f_j$ ist, was hier der Fall ist, da der einzige abhängige Fall für $i=j$ aus der SUmme rausgezogen würde, sodass $E(a b)=E(a)E(b)$ anwendbar ist.

Man beachte, dass der zweite Summand nur durch die Mittelung über mehrere komplett verschiedene Versuchsreihen Null wird. Betrachtet man jedoch nur eine Versuchsreihe, dann hat der zweite Summand auch ein Skalierverhalten in Abhängigkeit zu $N$. Da aber das Vorzeichen wechseln kann, muss man den Betrag betrachten:
\begin{align}
	\frac{1}{N} \sum\limits_{i=0}^N \left( f_i-\mu \right)
    \frac{1}{N} \sum\limits_{j=0,j\neq i}^N \left( f_j-\mu \right)
    =
	\frac{1}{N} \sum\limits_{i=0}^N
    \frac{1}{N} \sum\limits_{j=0,j\neq i}^N
    \left( f_i f_j - \mu \left( f_i + f_j \right) + \mu^2 \right)
    \\
    \label{eq:approxmuN}
    \approx
	\mu_N^2 - 2 \mu \mu_N + \mu^2
	\overset{ \mu_N \approx \mu - \sigma_{\mu_N} }{\approx}
	\mu^2 - 2 \sigma_{\mu_N} \mu + \sigma_{\mu_N}^2
	- 2\mu^2 + 2\mu \sigma_{\mu_N} + \mu^2
	= \sigma_{\mu_N}^2 = \frac{\sigma}{N}
\end{align}
Schritt \ref{eq:approxmuN} ist stark skizzenhaft und nicht mathematisch korrekt ausgeführt, wird aber gestützt durch empirische Auswertungen, vgl. Abb.\ref{fig:meanerrorsummand}. In der Abbildung sieht man, dass sowohl der erste Summand als auch der zweite invers proportional zu $N$ skaliert. Ein wichtiger Unterschied ist jedoch, dass der erste Summand immer positiv ist, während das Vorzeichen des zweiten Summanden oszilliert, wodurch er über die Mittelung mit $E(\cdot)$ gegen Null geht.

Interessant zu bemerken ist auch, dass der Graph der Standardvarianz des Mittelwertes $\sigma_{\mu_N}^2$ aufgetragen über die Anzahl an einbezogener Stichproben einer Zufallsbewegung ähnelt, anstatt stochastisch zu streuen. Dies wäre nicht der Fall, würde man für alle $N$ komplett neue Stichproben ziehen.

Weiterhin fällt auf, dass beide Summanden einer sehr glatten Geraden mit wenig Streuung folgen, während dies für $\sigma_{\mu_N}^2$ nicht der Fall ist. Dies zeigt, dass es durchaus zu einer Fehlerauslöschung durch den wegdiskutierten zweiten Summanden kommt. Dies beeinträchtigt jedoch nicht die Fehlerskalierung mit $\mathcal{O}\left( \frac{1}{N} \right)$.

\begin{figure}[H]
	\centering
	\begin{minipage}{0.7\linewidth}
		\includegraphics[width=\linewidth]{meanerror.pdf}
	\end{minipage}
	\caption{Darstellung der Standardvarianz des Mittelwertes $\sigma_{\mu_N}^2$ und der beiden in der Herleitung~\ref{eq:EabToEaEb-pre} auftretenden Summanden über die Anzahl einbezogener Stichproben. Hierbei ist zu beachten, dass beim Vergleich von den Werten für die Stichprobenanzahl von $N_1$ und $N_2$ die ersten $\min\left( N_1,N_2 \right)$ Stichproben identisch sind.}
	\label{fig:meanerrorsummand}
\end{figure}



%%%%%%%%%%%%%%%%%%%%%%%%%%%%%%%%%%%%%%%%%%%%%%%%%%%%%%%%%%%%%%%%%%%%%%%%%%%%%%%%
\chapter{Programmausdrücke}
%%%%%%%%%%%%%%%%%%%%%%%%%%%%%%%%%%%%%%%%%%%%%%%%%%%%%%%%%%%%%%%%%%%%%%%%%%%%%%%%

\lstinputlisting[language=bash, inputencoding=latin1, frame=single, caption={\lstinline!start_spark_slurm.sh! enthält Parameter für \lstinline!sbatch!, berechnet nötige Parameter für Spark aus denen von Slurm und startet dann einen Master und mehrere Worker-Prozesse startet}, label=lst:start_spark_slurm.sh]{../start_spark_slurm.sh}

\lstinputlisting[language=bash, inputencoding=latin1, frame=single, caption={\lstinline!startSpark.sh! welche eine Funktion definiert die einen den Sparkslurmjob aus Listing~\ref{lst:start_spark_slurm.sh} startet und Master-IP aus der Logdatei extrahiert}, label=lst:startSpark]{../startSpark.sh}

%\begin{lstlisting}[language=C++,firstnumber=10, caption={Die Trapezregel wie in Formel~\ref{eq:trapeze}}, label=lst:trapezregel]
%    for (uint64_t i=1; i<N-1; i++)
%    {
%    (!!!)
%     dx + f(b);
%\end{lstlisting}\end{minipage}\end{center}

%\begin{eqnarray}
%	\label{eq:n}
%	 a = b
%\end{eqnarray}
%
%\begin{figure}
%	\centering
%	\begin{minipage}{0.7 \linewidth}
%		\includegraphics[width=\linewidth]{Image.pdf}
%	\end{minipage}
%	\caption{captiontext}
%	\label{fig:RTn}
%\end{figure}


% add image(s) on one line without figure-environment, to place it exactly where you want
%\begin{minipage}{\linewidth}\begin{center}\bigskip
%	\captionsetup{type=figure}
%	\begin{minipage}{0.45 \linewidth}
%		\includegraphics[width=\linewidth]{ImageLeft.pdf}
%	\end{minipage}
%	\begin{minipage}{0.45 \linewidth}
%		\includegraphics[width=\linewidth]{ImageRight.pdf}
%	\end{minipage}
%	\captionof{figure}{...}
%	\label{fig:figlabel}
%\bigskip\end{center}\end{minipage}

\end{document}


% Typical problem words of me:
%  - versuchen
%  - zig
%  -
%  -
%  -

